\section{The Propositional Layer}
\label{sec:prop}

Before defining a Hoare logic we need a propositional logic to express
the pre- and post-conditions in.  To make explicit the connection to
forcing, we build this propositional logic as a Kripke-style forcing
semantics: a proposition is not a single truth value but a
\emph{family} of truth values indexed over a poset of
\emph{conditions}, and entailment is forcing at every condition.

The simplest forcing notion that supports step-indexed reasoning is
the natural numbers with the usual order.  Each natural number stands
for an ``observation depth'': condition $n$ records that we have
consumed $n$ program steps and committed to a partial view of the
computation.  Following the topos-of-trees and Iris convention, we
take \emph{smaller} naturals to be \emph{stronger} conditions: $0$ is
the weakest condition (we have observed nothing and are committed to
nothing), and larger $n$ are stronger (we are committed to a deeper
approximation).  Under this orientation, the forcing extension order
is the converse of $\le_\mathbb{N}$, and a forcing-style proposition
must be \emph{downward closed} in $\le_\mathbb{N}$: a proposition
forced at depth $n$ is forced at every shallower depth $m \le n$.

\begin{lstlisting}
Definition cond : Type := nat.

Record iProp : Type := MkProp {
  ihold : cond -> Prop;
  imono : forall n m, m <= n -> ihold n -> ihold m
}.

Notation "n ⊩ φ" := (ihold φ n).
\end{lstlisting}

The notation $n \Vdash \varphi$ matches the classical forcing notation,
and it is more than analogical: an \texttt{iProp}, taken with its
\texttt{imono} field, is exactly a presheaf on $(\omega, \le)$, hence
an object of the topos of trees \cite{birkedal-et-al-2012}, which
is exactly the Kripke model of intuitionistic logic with worlds in
$\omega$, and hence (Section~\ref{sec:topos}) the forcing-style
semantics with conditions in $\omega$.

Entailment is forcing at every condition:
\begin{lstlisting}
Definition ientails (φ ψ : iProp) : Prop :=
  forall n, n ⊩ φ -> n ⊩ ψ.

Notation "φ ⊢ ψ" := (ientails φ ψ).
\end{lstlisting}

\paragraph{Connectives.}
The Heyting-algebra connectives have the standard Kripke
interpretation.  Truth and falsity are the constant predicates.
Conjunction and disjunction are pointwise.  Implication is the
Kripke implication, quantified over all stronger conditions:

\begin{lstlisting}
Definition iAnd (φ ψ : iProp) : iProp :=
  MkProp (fun n => n ⊩ φ /\ n ⊩ ψ) (iAnd_mono φ ψ).

Definition iImpl (φ ψ : iProp) : iProp :=
  MkProp (fun n => forall m, m <= n -> m ⊩ φ -> m ⊩ ψ)
         (iImpl_mono φ ψ).
\end{lstlisting}

The quantification ``$\forall m \le n$'' in \texttt{iImpl} is the
characteristic move of intuitionistic Kripke semantics. Without it,
the framework would collapse to classical logic with a step index.
The monotonicity proofs \texttt{iAnd\_mono} and \texttt{iImpl\_mono}
are routine. In each case the proof simply transports the membership
witness through the order.

\paragraph{The later modality.}
The genuinely new ingredient over the standard Kripke interpretation
of intuitionistic logic is the \emph{later} modality $\triangleright$.
It is the operator that says ``to talk about a proposition here, you
must drop one level''.  In our $(\omega, \le)$ orientation, $\triangleright \varphi$ at depth $n$ holds if $n = 0$ (we are at the weakest condition and the modality is vacuously satisfied) or if $\varphi$ holds at the shallower depth $n - 1$:

\begin{lstlisting}
Definition iLater (φ : iProp) : iProp :=
  MkProp (fun n => match n with
                   | 0 => True
                   | S k => k ⊩ φ
                   end) (iLater_mono_helper φ).

Notation "▷ φ" := (iLater φ).
\end{lstlisting}

Two readings of this operator are worth stating explicitly.  In
step-indexed-logic terms, $\triangleright \varphi$ is ``$\varphi$ when
we look at the computation one step later'', and the shift is what
guards self-referential definitions against circularity.  In
forcing terms, $\triangleright$ is the explicit level-descent
operator: at level $n+1$ we may consult level $n$.  This is precisely
the ramification of Cohen's original construction (Section~\ref{sec:related}, ``Forcing in set theory''): the recursive
definition of forcing at level $\alpha$ may invoke forcing only at
strictly lower levels, and $\triangleright$ is the syntactic device
that makes the descent explicit.

\paragraph{Heyting algebra rules.}
With the connectives in hand, the standard intuitionistic
introduction and elimination rules become direct consequences of
their Kripke definitions.  We mechanize the full set:

\begin{lstlisting}
Lemma iAnd_intro φ ψ χ : φ ⊢ ψ -> φ ⊢ χ -> φ ⊢ ψ ∧ χ.
Lemma iAnd_proj_l φ ψ : φ ∧ ψ ⊢ φ.
Lemma iOr_intro_l φ ψ : φ ⊢ φ ∨ ψ.
Lemma iOr_elim φ ψ χ : φ ⊢ χ -> ψ ⊢ χ -> φ ∨ ψ ⊢ χ.
Lemma iImpl_intro φ ψ χ : φ ∧ ψ ⊢ χ -> φ ⊢ (ψ → χ).
Lemma iImpl_elim φ ψ : (φ → ψ) ∧ φ ⊢ ψ.
\end{lstlisting}

\noindent
Each proof simply moves the assumed monotonicity through the
relevant connective. The work is small but worth doing
mechanically because it confirms that the framework is genuinely a
Heyting algebra without any informal appeal to topos-theoretic
abstractions.

The basic laws of the later modality follow the same pattern.  We
prove that $\triangleright$ is monotone in the entailment order,
that any proposition entails its own later
($\varphi \vdash \triangleright \varphi$), that $\triangleright$
distributes over $\wedge$ in both directions, and that
$\triangleright \top$ is interderivable with $\top$.

\paragraph{Löb's rule.}
The technical heart of the propositional layer is Löb's rule:
\begin{theorem}[Löb's rule]
For every \texttt{iProp} $\varphi$, $(\triangleright \varphi \to
\varphi) \vdash \varphi$.
\end{theorem}

\noindent
We give the rule the standard step-indexed proof, by induction on
$n$:

\begin{lstlisting}
Lemma iLob φ : (▷ φ → φ) ⊢ φ.
Proof.
  intros n Hn.
  induction n as [|n IH].
  - apply (Hn 0); [lia | exact I].
  - apply (Hn (S n)); [lia | simpl].
    apply IH.
    intros m Hm Hlate.
    apply (Hn m); [lia | exact Hlate].
Qed.
\end{lstlisting}

It is worth noticing what this proof actually uses.  The base case
relies on the fact that $\triangleright \varphi$ at $n=0$ is
\emph{vacuously} \texttt{True}, while the inductive step uses the
hypothesis at $n+1$ together with the induction hypothesis at $n$.
In other words, Löb's rule is, at the meta level, well-founded
induction on the forcing order $\omega$.  This observation will
return in Section~\ref{sec:abstract}: we will see there that in the
abstract framework Löb's rule \emph{is} \texttt{well\_founded\_ind}
on the strict refinement relation.

\paragraph{Step-indexing is not collapsing.}
A useful sanity check is that the later modality changes the
propositional theory: it is not the case that
$\triangleright \bot$ entails $\bot$.  At $n = 0$, the proposition
$\triangleright \bot$ is \texttt{True} by definition, whereas $\bot$
is \texttt{False}, so the entailment fails at $n = 0$:

\begin{lstlisting}
Lemma iLater_False_distinct : ~ (▷ ⊥ ⊢ ⊥).
Proof. intros H. apply (H 0); simpl; exact I. Qed.
\end{lstlisting}

\noindent
This is the canonical witness that step-indexing genuinely extends
the propositional theory, rather than collapsing to ordinary
intuitionistic logic.  We will revisit it in
Section~\ref{sec:lob-essential} to make the converse point: if we
\emph{drop} the well-foundedness of the forcing order, then
$\triangleright \bot$ becomes equivalent to $\bot$ in a way that
breaks Löb's rule entirely.

